\documentclass[11pt]{article}
\usepackage[margin=1in]{geometry}
\usepackage{amsmath,amssymb,amsthm}
\usepackage{hyperref}
\usepackage{enumitem}
\usepackage{xurl}

\newtheorem{theorem}{Theorem}
\newtheorem{lemma}[theorem]{Lemma}
\theoremstyle{definition}
\newtheorem{definition}[theorem]{Definition}
\theoremstyle{remark}
\newtheorem{remark}[theorem]{Remark}

\let\oldus\_
\renewcommand{\_}{\oldus\allowbreak}

\title{A disproof of Conjecture 00000007661\\[2pt]
\large Q-geometric ratios of q-hypergeometric integer series}
\date{}

\begin{document}
\sloppy
\maketitle

\section{The conjecture}

Conjecture 00000007661 reads:

\begin{quote}
\emph{Definition:} Let $0<|q|<1$; a power series $F(z)=\sum a_n z^n$ is of \emph{q-hypergeometric type}
if there are rational functions $R(z),S(z)$ with $F(qz)=R(z)F(z)+S(z)$; the coefficient ratios
$a_{n+1}/a_n$ are the \emph{Q-geometric ratios}. \emph{Conjecture:} $F(z)\in\mathbb Z[[z]]$ is of
q-hypergeometric type with convergence radius 1 if and only if its Q-geometric ratio sequence converges to
1 at rate $O(q^n)$, and this rate cannot be improved to $o(q^n)$ unless $F$ is a polynomial.
\end{quote}

Fix $q\in\mathbb C$ with $0<|q|<1$ (the Lean proof works for every such $q$, so it refutes both the
reading ``for every such $q$'' and the reading ``for some such $q$''). Write
$\mathrm{L}(F)$ for ``$F$ is of q-hypergeometric type and has radius of convergence $1$'',
$\mathrm{O}(F)$ for ``$a_{n+1}/a_n\to1$ and $a_{n+1}/a_n-1=O(q^n)$'', and
$\mathrm{o}(F)$ for ``$a_{n+1}/a_n-1=o(q^n)$''. We refute the following three readings:
\begin{enumerate}[label=(\Alph*)]
  \item the equivalence $\mathrm{L}(F)\iff\mathrm{O}(F)$ for all $F\in\mathbb Z[[z]]$;
  \item the sharpness clause $\mathrm{o}(F)\Rightarrow F$ is a polynomial, even restricted to $F$ with
        $\mathrm{L}(F)$;
  \item the equivalence $\mathrm{L}(F)\iff\bigl(\mathrm{O}(F)\wedge(\mathrm{o}(F)\Rightarrow F\text{ is a
        polynomial})\bigr)$, in which the sharpness clause is part of the right-hand side.
\end{enumerate}
The statement read as (A) and (B) together is false because (A) is false (and also because (B) is false).

\section{Definitions and conventions}

\begin{itemize}
  \item $F=\sum a_nz^n\in\mathbb Z[[z]]$ is evaluated as the complex sum $F(z)=\sum_{n\ge0}a_nz^n$
        (Lean: \texttt{toFun}, a \texttt{tsum}); the radius of convergence is the radius of the complex
        power series $\sum a_nz^n$ (Lean: \texttt{radius}, the radius of
        \texttt{FormalMultilinearSeries.ofScalars}).
  \item A rational function is $A/B$ with $A,B\in\mathbb C[z]$, $B\ne0$. Analytic reading
        (\texttt{QHypergeometricType}): there are $A,B,C,D\in\mathbb C[z]$ with $B,D\ne0$ such that
        $F(qz)=\frac{A(z)}{B(z)}F(z)+\frac{C(z)}{D(z)}$ for every $|z|<1$ with $B(z)\ne0\ne D(z)$.
        Formal reading (\texttt{QHypergeometricTypeFormal}): $BD\cdot F(qz)=AD\cdot F+BC$ in
        $\mathbb C[[z]]$, where $F(qz)$ is the series $\sum q^na_nz^n$ (Mathlib \texttt{rescale q}).
  \item The Q-geometric ratio is $a_{n+1}/a_n$, computed in $\mathbb R$ (Lean: \texttt{ratio}). Both
        examples below have all $a_n\neq0$, so no convention for division by zero is used.
  \item $O(q^n)$ and $o(q^n)$ are Mathlib's \texttt{IsBigO} and \texttt{IsLittleO} along
        $n\to\infty$; they compare $|a_{n+1}/a_n-1|$ with $|q|^n$.
  \item ``$F$ is a polynomial'' means $F$ is the image of some $p\in\mathbb Z[z]$ (\texttt{IsPoly}).
\end{itemize}

\section{The counterexamples}

Let $F_1=\sum_{n\ge0}z^n=1/(1-z)$ and $F_2=\sum_{n\ge0}(n+1)z^n=1/(1-z)^2$, both in $\mathbb Z[[z]]$.

\begin{lemma}\label{lem:qhyp}
For $|z|<1$, $F_1(z)=(1-z)^{-1}$ and $F_2(z)=(1-z)^{-2}$. Hence
$F_1(qz)=\frac{1-z}{1-qz}F_1(z)$ and $F_2(qz)=\frac{(1-z)^2}{(1-qz)^2}F_2(z)$ on the unit disk, and
both are of q-hypergeometric type ($S=0$) in the analytic and the formal sense.
\end{lemma}
\begin{proof}
The geometric series gives $F_1$; $\sum nz^n=z/(1-z)^2$ (Mathlib
\texttt{hasSum\_coe\_mul\_geometric\_of\_norm\_lt\_one}) plus $\sum z^n$ gives $F_2$. Since $|qz|<1$,
the same formulas hold at $qz$, and $1-z\ne0$. Formally, $(1-z)F_1=1$ and $(1-z)^2F_2=1$ in
$\mathbb C[[z]]$ (Mathlib \texttt{mk\_one\_mul\_one\_sub\_eq\_one},
\texttt{mk\_add\_choose\_mul\_one\_sub\_pow\_eq\_one}); applying the ring homomorphism
$z\mapsto qz$ gives $(1-qz)F_1(qz)=1$ and $(1-qz)^2F_2(qz)=1$, so both sides of the cleared identities equal
$1$. (Lean: \texttt{qHyp\_F1}, \texttt{qHyp\_F2}, \texttt{qHypFormal\_F1}, \texttt{qHypFormal\_F2}.)
\end{proof}

\begin{lemma}
Both $F_1$ and $F_2$ have radius of convergence $1$.
\end{lemma}
\begin{proof}
$|a_{n+1}|/|a_n|$ equals $1$ for $F_1$ and $(n+2)/(n+1)=1+1/(n+1)\to1$ for $F_2$; by the ratio test
(Mathlib \texttt{ofScalars\_radius\_eq\_inv\_of\_tendsto}) the radius is $1$.
(Lean: \texttt{radius\_F1}, \texttt{radius\_F2}.)
\end{proof}

\begin{lemma}
(i) For $F_2$, $a_{n+1}/a_n-1=1/(n+1)$, which is not $O(q^n)$.
(ii) For $F_1$, $a_{n+1}/a_n=1$, so $\mathrm{O}(F_1)$ and $\mathrm{o}(F_1)$ hold; $F_1$ is not a polynomial.
\end{lemma}
\begin{proof}
(i) If $1/(n+1)\le c|q|^n$ for all large $n$, then $1\le c(n+1)|q|^n$, but $(n+1)|q|^n\to0$ (Mathlib
\texttt{tendsto\_self\_mul\_const\_pow\_of\_lt\_one}). (ii) The zero sequence is $o$ of anything. If
$F_1=p$ for a polynomial $p$, the coefficient of $z^{\deg p+1}$ would be both $1$ and $0$.
(Lean: \texttt{not\_rateBigO\_F2}, \texttt{rateBigO\_F1}, \texttt{rateLittleO\_F1}, \texttt{not\_isPoly\_F1}.)
\end{proof}

\begin{theorem}
For every $q\in\mathbb C$ with $0<|q|<1$, readings (A), (B), (C) are all false, with either reading of
``q-hypergeometric type''.
\end{theorem}
\begin{proof}
(A): $\mathrm{L}(F_2)$ holds but $\mathrm{O}(F_2)$ fails. (B): $\mathrm{L}(F_1)$ and $\mathrm{o}(F_1)$ hold
but $F_1$ is not a polynomial. (C): $\mathrm{L}(F_1)$ holds but the right-hand side fails at $F_1$, since
$\mathrm{o}(F_1)$ holds and $F_1$ is not a polynomial.
\end{proof}

\section{Lean formalization}

File \texttt{lean/Conjecture7661/Basic.lean} (namespace \texttt{C7661}, Lean 4 + Mathlib, only
\texttt{import Mathlib}). Main theorems:
\begin{itemize}
  \item \texttt{conjecture\_false (q : C) (\_hq0 : 0 < norm q) (hq1 : norm q < 1)} (ASCII rendering): the facts about $F_2$ and
        $F_1$ above, and the negations of (A), (B) (restricted to $\mathrm{L}$) and (C), with the analytic
        definition \texttt{QHypergeometricType}.
  \item \texttt{conjecture\_false\_formal}: the negations of (A), (B), (C) with the formal definition
        \texttt{QHypergeometricTypeFormal}.
\end{itemize}
Axiom audit: \texttt{\#print axioms} for both theorems reports only \texttt{propext},
\texttt{Classical.choice}, \texttt{Quot.sound}. There is no \texttt{sorry} and no \texttt{native\_decide}.

\begin{remark}
We do not refute any reading in which the right-hand side does not imply $a_{n+1}/a_n-1=O(|q|^n)$ and the
sharpness clause is dropped; we know of no such reading of the text.
\end{remark}

\end{document}
